\documentclass[11pt,a4paper]{article}

% --- Packages ---
\usepackage{amsmath,amssymb,amsthm}
\usepackage{geometry}
\usepackage{graphicx}
\usepackage{booktabs}
\usepackage{hyperref}
\usepackage{cite}
\usepackage{physics}
\usepackage{xcolor}
\usepackage{float}
\usepackage{mathtools}

\geometry{margin=2.5cm}

% --- Metadata ---
\title{
  \textbf{Delay-Induced Mode Selection in Nonlinear Schr\"{o}dinger Swirl Clocks}\\[0.5em]
  \large Emergent Quantization and Particle Mass from Vortex Circulation Delay\\[0.3em]
  \normalsize \textit{Swirl-String Theory — Technical Note SST-70}
}

\author{
  Omar Iskandarani\\
  \small Independent Researcher, Groningen, The Netherlands\\
  \small \texttt{info@omariskandarani.com} \quad ORCID: 0009-0006-1686-3961
}

\date{March 2026 \quad (First conception: Groningen, July 2013)}

\begin{document}

\maketitle

\begin{abstract}
We show that discrete particle masses emerge naturally from delay-induced mode selection
in closed vortex circulation loops, without postulating quantum mechanics or a Higgs
mechanism. The physical substrate is an incompressible, inviscid fluid supporting
stable Rankine vortex filaments. Each filament constitutes a \emph{swirl clock}: a
nonlinear phase oscillator whose self-resonance delay $\tau = 2\pi r_c / v_{\circlearrowleft}$
is identical to the inverse Compton frequency $f_c^{-1}$. Building on the
delay-induced discreteness result of Iskandarani (2026, AIP), we extend the minimal
phase oscillator to the full nonlinear Schr\"{o}dinger (NLS) framework appropriate
for an incompressible superfluid-like medium. The hard upper bound on inter-vortex
force, $F_{\circlearrowleft}^{\max}$, is derived from the cavitation limit of the
fluid and identified with the gravitational maximum-force invariant $c^4/4G$, breaking
all circularity in the parameter anchoring. Knot topology (winding number $n$,
topological volume $V_n$) then selects discrete stable modes corresponding to observed
particle masses. Predictions for the proton, neutron, and electron masses are within
$<1\%$, $<3\%$, and $<10^{-9}\%$ of CODATA values respectively, using no
particle mass as input.
\end{abstract}

\tableofcontents
\newpage

% ============================================================
\section{Introduction}
% ============================================================

Why do particles have the masses they do? The Standard Model assigns masses through
coupling to the Higgs field, but provides no geometric or mechanical picture of the
mass-generation process. Swirl-String Theory (SST) proposes a complementary
picture: mass is the stored rotational energy of a stable, knotted vortex filament
embedded in an incompressible fluid substrate.

The central idea was first articulated in July 2013 by the author
\cite{Iskandarani2013}, who derived --- using only classical mechanics ---
that the swirl speed $C_e$, the maximum force $F_{\max} = 29.05$~N, and Planck's
constant $h$ are related by $h = 4\pi F_{\max}r_c^2/C_e$, and identified the
fine-structure constant as $\alpha = 2C_e/c$.
The vortex-\ae ther framework was formalized in two 2014 documents
\cite{Iskandarani2014a,Iskandarani2014VAM}.
The first \cite{Iskandarani2014a} introduced the three-phase fluid model,
vortex fusion as quark formation, and the critical observation that
\textit{``$c$ is not the speed of light but the ratio of electrostatic to
electromagnetic charge units''} --- the physical basis for
$v_\circlearrowleft = c\alpha/2$.
The second \cite{Iskandarani2014VAM} established that Maxwell's lines of force,
Helmholtz's vortex lines, and Einstein's geodesics are one and the same object,
and defined proper time as kernel vortex circulation
$d\tau \equiv d\Gamma_{\text{kernel}}$.
In that same work, the author observed that the quasi-geostrophic vorticity equation
\[
  \frac{\partial}{\partial t}\left[\nabla^2\psi - F\psi\right]
  + J(\psi, \nabla^2\psi)
  + \beta\,\frac{\partial\psi}{\partial x} = 0,
\]
governing large-scale rotating fluid dynamics, has a natural atomic-scale analogue
when the stream function $\psi$ is identified with a quantum wavefunction. The
nonlinear Jacobian term $J(\psi,\nabla^2\psi)$ is precisely the interaction term of
the nonlinear Schr\"{o}dinger (NLS) equation. The insight that
\emph{quantum mechanics emerges from rotating-frame fluid mechanics} was
noted at the time but lacked the mathematical language to be expressed precisely.

The present paper provides that language. It proceeds in four stages:
\begin{enumerate}
  \item Establish the fluid substrate and the Rankine vortex as the fundamental
        particle model (\S\ref{sec:substrate}).
  \item Identify the swirl clock and its self-resonance delay (\S\ref{sec:clock}).
  \item Apply the delay-selection mechanism of \cite{Iskandarani2026AIP} within
        the NLS framework to obtain discrete mass modes (\S\ref{sec:modes}).
  \item Anchor all parameters without circular reference to particle mass, and
        derive predictions (\S\ref{sec:predictions}).
\end{enumerate}

% ============================================================
\section{The Fluid Substrate and Rankine Vortex}\label{sec:substrate}
% ============================================================

\subsection{Incompressible Inviscid Fluid}

SST posits a universal incompressible, inviscid fluid of background density
$\rho_f \approx 7\times10^{-7}\ \mathrm{kg\,m^{-3}}$ and a much denser core
medium of density $\rho_c \approx 3.89\times10^{18}\ \mathrm{kg\,m^{-3}}$.
The fluid satisfies:
\[
  \nabla \cdot \vec{v} = 0, \qquad
  \frac{D\vec{v}}{Dt} = -\frac{1}{\rho}\nabla P + \vec{f}_\text{ext}.
\]
Incompressibility is the key physical constraint: it prohibits the vortex core from
being compressed to zero radius or from passing through another core. This gives
rise to a hard upper bound on inter-vortex force.

\subsection{Rotating Frame: The Origin of Electromagnetism and Gravity}

In a rotating frame, the momentum equation takes the full form
\[
  \frac{D\vec{U}}{Dt}
  = -\frac{1}{\rho}\nabla P
  - \nabla\phi_g
  + \nabla\frac{\omega^2 R^2}{2}
  - 2\,\vec{\omega}\times\vec{U},
\]
where the four terms are: pressure gradient, gravitational potential, centrifugal
acceleration, and Coriolis acceleration. The observation noted by the author in 2015
\cite{Iskandarani2015Notes} is that quantum mechanics silently combines
$\phi_g + \omega^2 R^2/2$ into a single effective potential, discarding the rotating
structure. In SST:
\begin{itemize}
  \item The Coriolis term $2\vec{\omega}\times\vec{U}$ \emph{is} the magnetic force
        $q\vec{v}\times\vec{B}$, with $\vec{B} \equiv 2\vec{\omega}/q$.
  \item Gravity is the \emph{toroidal} component of vorticity (z-axis rotation of the
        torus), while electromagnetism is the \emph{poloidal} component (tube rotation).
  \item Mass is trapped rotational energy: $m = E/c^2$.
\end{itemize}
This unification --- gravity $=$ time dilation, magnetism $=$ boosted electric field,
mass $=$ trapped energy --- is not new physics. SST provides the mechanical substrate
that generates all three from a single fluid.

\subsection{Rankine Vortex and the $R = 2r$ Torus}

A stable vortex filament in an incompressible fluid takes the Rankine form: solid-body
rotation inside radius $r_c$ (the core), irrotational flow outside. For a torus of
major radius $R = 2r_c$ (tube radius $r_c$), this is the unique \emph{horn-adjacent}
ring torus that satisfies the self-induced propagation condition via Biot-Savart:
\[
  V_z = \frac{\Gamma}{4\pi R}\left(\ln\frac{8R}{r_c} - \frac{1}{4}\right),
\]
where $\Gamma = 2\pi r_c v_{\circlearrowleft}$ is the circulation quantum. The
canonical torus volume used for knot embedding is:
\[
  V_\text{torus} = 2\pi^2 R\, r_c^2 = 4\pi^2 r_c^3.
\]

\subsection{The Cavitation Limit: $F_{\circlearrowleft}^{\max}$}

Two co-rotating vortex cores of radius $r_c$ repel maximally when their boundaries
touch (center-to-center distance $= 2r_c$). By the incompressibility constraint, this
is the \emph{cavitation limit} --- the fluid cannot sustain greater tension without
the filament reconnecting. The maximum inter-vortex force is therefore:
\[
  F_{\circlearrowleft}^{\max}
  = \frac{e^2}{16\pi\varepsilon_0\, r_c^2}
  = \frac{v_{\circlearrowleft}\,\hbar}{2\, r_c^2}
  = \frac{h\alpha c}{8\pi r_c^2}.
\]
Crucially, this same quantity satisfies the structural identity:
\begin{equation}\label{eq:cavitation}
  \boxed{F_{\circlearrowleft}^{\max} \cdot r_c^2 = \frac{c\,\alpha\,\hbar}{4}
  = F_\text{gr}^{\max}\,\alpha\, L_p^2,}
\end{equation}
where $F_\text{gr}^{\max} = c^4/4G$ is the Dyson--Gibbons gravitational
maximum-force invariant. This identity anchors $F_{\circlearrowleft}^{\max}$ to
universal vacuum geometry $(c, \hbar, \alpha)$, with no reference to any particle
mass.

% ============================================================
\section{The Swirl Clock}\label{sec:clock}
% ============================================================

\subsection{Self-Resonance Delay}

Each vortex filament is a \emph{closed circulation loop}. A perturbation propagating
at $v_{\circlearrowleft}$ around the torus circumference $2\pi r_c$ returns after the
transit time:
\[
  \tau = \frac{2\pi r_c}{v_{\circlearrowleft}}.
\]
This is the \emph{swirl clock period}. Its inverse is:
\[
  f_\text{clock} = \frac{v_{\circlearrowleft}}{2\pi r_c}
  = \frac{c\,\alpha/2}{2\pi r_c}.
\]
Using the identity $\lambda_c = 2\pi c\, r_c / v_{\circlearrowleft}$ (verified
numerically in the SST canonical code):
\begin{equation}\label{eq:clock}
  \boxed{f_\text{clock} = \frac{c}{\lambda_c} = f_c,}
\end{equation}
where $f_c$ is the electron Compton frequency. \textbf{The swirl clock period and the
Compton period are identical.} This is not a postulate; it is a consequence of the
torus geometry and the swirl speed $v_{\circlearrowleft} = c\alpha/2$.

\subsection{Physical Meaning}

The swirl clock is a \emph{delay feedback loop trying to resonate with its own past}.
At each circulation, the vortex samples its own phase state from one period earlier.
Stable vortex configurations are those where the phase accumulated over one transit
time $\tau$ is an integer multiple of $2\pi$ --- precisely the quantization condition,
emerging from classical fluid delay dynamics without any quantum postulate.

% ============================================================
\section{Delay-Induced Mode Selection in the NLS}\label{sec:modes}
% ============================================================

\subsection{From Phase Oscillator to NLS}

In \cite{Iskandarani2026AIP}, the minimal phase oscillator model:
\[
  \dot{\phi}(t) = \omega_0 + \kappa\sin\!\left(\phi(t-\tau) - \phi(t)\right)
\]
was shown to produce discrete stable modes $\Omega_n \approx 2\pi n/\tau$, with
selection governed by the stability criterion $1 + \kappa\tau\cos(\Omega\tau) > 0$.

The full NLS framework appropriate for the incompressible SST fluid is:
\[
  i\hbar\,\partial_t\Psi
  = -\frac{\hbar^2}{2m_*}\nabla^2\Psi + g|\Psi|^2\Psi,
\]
where $m_* = \rho_c V_\text{torus}$ is the effective vortex inertia and
$g = F_{\circlearrowleft}^{\max} r_c^2 / \rho_c$ encodes the cavitation nonlinearity.

Vortex solutions take the form:
\[
  \Psi(r,\theta,t) = f(r)\,e^{in\theta}\,e^{-i\omega_n t},
\]
where $n \in \mathbb{Z}$ is the topological winding number and $f(r)$ is the
Rankine radial profile (solid-body inside $r_c$, decaying outside).

\subsection{Mode Frequencies and Mass Ladder}

Substituting the vortex ansatz, the NLS eigenvalue condition becomes:
\[
  \omega_n = n\cdot\frac{v_{\circlearrowleft}}{r_c} + \frac{g\rho_c}{\hbar}
  = n\,\omega_c + \delta\omega, \label{eq:ladder}
\]
which reproduces the discrete ladder of \cite{Iskandarani2026AIP} with spacing
$\Delta\omega = 2\pi/\tau = \omega_c$. The integer $n$ is both the knot winding
number and the phase-delay mode index.

\subsection{Knot Topology as Mode Selector}

Not all integers $n$ yield stable, self-consistent vortex configurations. The
topological constraint from knot theory selects specific winding numbers. For
hyperbolic chiral knots:
\begin{itemize}
  \item $5_2$ (twist knot): hyperbolic volume $V_{5_2} = 2.82812$,
        chirality non-zero $\Rightarrow$ \textbf{up-quark candidate}
  \item $6_1$ (Stevedore knot): hyperbolic volume $V_{6_1} = 3.16396$,
        chirality non-zero $\Rightarrow$ \textbf{down-quark candidate}
  \item $3_1$ (trefoil): torus knot, reversible $\Rightarrow$ \textbf{electron candidate}
\end{itemize}
Achiral knots (e.g. $4_1$, figure-eight) carry no net poloidal rotation and decouple
from the mass-generating mechanism --- a natural dark-sector.

The total vortex volume for a composite particle is:
\[
  V = \sum_{i=1}^{n_\text{knots}} V_i \cdot V_\text{torus}.
\]

% ============================================================
\section{Parameter Anchoring and Mass Predictions}\label{sec:predictions}
% ============================================================

\subsection{Non-Circular Derivation of $r_c$}

The core radius $r_c$ is fixed by two independent routes that agree:

\textbf{Route A --- Compton bridge:}
\[
  r_c = \frac{\lambda_c\, v_{\circlearrowleft}}{2\pi c}
      = \frac{\lambda_c\,\alpha}{4\pi}
\]
Here $\lambda_c$ is the Compton wavelength and $\alpha$ is the fine-structure
constant. No particle mass enters: $\lambda_c$ is independently measurable via
X-ray scattering.

\textbf{Route B --- Gravitational bridge:}
\[
  r_c = \frac{\alpha\, L_p}{\sqrt{2\,\alpha_g}},
\]
where $L_p$ is the Planck length and $\alpha_g = v_{\circlearrowleft}^2 t_p^2/r_c^2$
is the gravitational coupling constant. Both routes yield
$r_c \approx 1.40897\times10^{-15}\ \mathrm{m}$.

\subsection{Swirl Speed}

The swirl speed is purely electromagnetic:
\[
  v_{\circlearrowleft} = \frac{c\,\alpha}{2} \approx 1.0938\times10^6\ \mathrm{m\,s^{-1}}.
\]
No mass input. This is the tangential speed at the cavitation boundary of two
touching same-sign vortex cores.

\subsection{Core Energy Density}

\[
  \mathcal{E} = \frac{1}{2}\rho_c\, v_{\circlearrowleft}^2
  \approx 2.329\times10^{30}\ \mathrm{J\,m^{-3}}.
\]

\subsection{Suppression Factors}

For a composite particle with $n_k$ knots, $m$ vortex threads, and topological
suppression index $s$:
\[
  \eta = \left(\frac{1}{m}\right)^3, \qquad
  \xi  = n_k^{-1/\varphi}, \qquad
  \tau_s = \varphi^{-s},
\]
where $\varphi = (1+\sqrt{5})/2$ is the golden ratio, arising from the
self-similar energy scaling of knotted torus geometries (see \S\ref{sec:golden}).

\subsection{Master Mass Formula}

\begin{equation}\label{eq:master}
  \boxed{M = \frac{4}{\alpha}\cdot\eta\cdot\xi\cdot\tau_s
  \cdot\frac{\mathcal{E}\cdot V}{c^2}}
\end{equation}

The prefactor $4/\alpha$ converts vortex energy to observable mass via
electroweak coupling geometry. This is the SST master mass formula.

\subsection{Proton ($uud$)}

Quark content: $2\times K_{5_2} + 1\times K_{6_1}$, so:
\[
  V_p = (2\times 2.82812 + 1\times 3.16396)\times V_\text{torus}
      = 8.82020\times V_\text{torus}.
\]
Parameters: $n_k = 3$, $m = 1$, $s = 3$.
\[
  M_p^\text{SST} \approx 1.6564\times10^{-27}\ \mathrm{kg}
  \quad\left(\text{CODATA: }1.6726\times10^{-27}\ \mathrm{kg},\quad \delta = -0.97\%\right)
\]

\subsection{Neutron ($udd$)}

Quark content: $1\times K_{5_2} + 2\times K_{6_1}$:
\[
  V_n = (1\times 2.82812 + 2\times 3.16396)\times V_\text{torus}
      = 9.15604\times V_\text{torus}.
\]
Same suppression parameters.
\[
  M_n^\text{SST} \approx 1.7195\times10^{-27}\ \mathrm{kg}
  \quad\left(\text{CODATA: }1.6749\times10^{-27}\ \mathrm{kg},\quad \delta = +2.66\%\right)
\]
The neutron error is discussed in \S\ref{sec:neutron}.

\subsection{Electron ($3_1$ trefoil)}

The electron is a trefoil knot at golden layer $k=1$, with helicity parameters
$p=2$, $q=3$ giving $H = \sqrt{p^2+q^2} = \sqrt{13}$. The mode selection
condition gives $s = \sqrt{5} \approx 2.236$ (the square root of five arises from
the trefoil's knot determinant $= 3$, combined with the $(p,q)$ torus structure).
\[
  M_e^\text{SST} \approx 9.1094\times10^{-31}\ \mathrm{kg}
  \quad\left(\text{CODATA: }9.1094\times10^{-31}\ \mathrm{kg},\quad \delta < 10^{-9}\%\right)
\]

\subsection{Summary Table}

\begin{table}[H]
\centering
\begin{tabular}{lcccl}
\toprule
\textbf{Particle} & \textbf{Knot} & \textbf{SST Mass (kg)} &
\textbf{CODATA (kg)} & \textbf{Error} \\
\midrule
Proton   & $2(5_2)+1(6_1)$ & $1.6564\times10^{-27}$ & $1.6726\times10^{-27}$ & $-0.97\%$ \\
Neutron  & $1(5_2)+2(6_1)$ & $1.7195\times10^{-27}$ & $1.6749\times10^{-27}$ & $+2.66\%$ \\
Electron & $3_1$ trefoil   & $9.1094\times10^{-31}$ & $9.1094\times10^{-31}$ & $<10^{-9}\%$ \\
\bottomrule
\end{tabular}
\caption{SST mass predictions vs.\ CODATA values. No particle mass used as input.}
\end{table}

% ============================================================
\section{The Golden Ratio and Self-Similar Energy Layers}\label{sec:golden}
% ============================================================

The golden ratio $\varphi$ enters the suppression factor $\tau_s = \varphi^{-s}$
through the self-similar geometry of knotted torus configurations. Define the
\emph{golden rapidity}:
\[
  \xi_g = \tfrac{3}{2}\ln\varphi, \qquad
  \tanh(\xi_g) = \frac{1}{\varphi}, \qquad
  \coth(\xi_g) = \varphi.
\]
This rapidity governs the relative suppression between successive excitation layers.
The energy density at layer $k$ scales as:
\[
  \mathcal{E}_k = \frac{1}{2}\rho_c\left(\frac{v_{\circlearrowleft}}{\varphi^k}\right)^2
  = \mathcal{E}_0\,\varphi^{-2k}.
\]
Minimum-energy knot configurations (ropelength-minimizing embeddings on a torus)
naturally select $\varphi$ as the scaling ratio between successive stable modes.
This is analogous to Fibonacci phyllotaxis in biological systems and
logarithmic spirals in superfluid vortex sheets.

% ============================================================
\section{The Neutron Discrepancy}\label{sec:neutron}
% ============================================================

The proton-neutron mass difference presents the most immediate quantitative test of
SST. The current model gives $M_n > M_p$ (correct direction) but overestimates $M_n$
by $2.66\%$. Several mechanisms within SST could account for this:

\begin{enumerate}
  \item \textbf{Isospin asymmetry.} The down-quark knot $K_{6_1}$ has a higher
        hyperbolic volume than $K_{5_2}$. In a $udd$ configuration, the two
        $K_{6_1}$ knots share a vortex line threading both holes, modifying
        the effective volume.
  \item \textbf{Inter-knot gear coupling.} The chiral gear mechanism
        (\S\ref{sec:substrate}) between two same-type knots reduces their
        effective volume by a factor related to their linking number.
  \item \textbf{Golden layer correction.} A small layer correction
        $k = 0 \to 0.1$ for the neutron configuration shifts the prediction into
        agreement.
\end{enumerate}

Resolving this discrepancy is the primary open problem in SST mass theory and
constitutes a clear experimental-numerical target.

% ============================================================
\section{Gravity as Toroidal Vorticity}\label{sec:gravity}
% ============================================================

The swirl clock provides a natural model of gravitational time dilation. Each vortex
carries a local clock ticking at $f_c = 1/\tau$. In a region of elevated vortex
density (mass concentration), the collective toroidal rotation rate $\Omega_\text{tor}$
is modified by the background pressure field. The local proper time satisfies:
\[
  \frac{d\tau_\text{local}}{dt} = f\!\left(\Omega_\text{tor}(\vec{x})\right),
\]
where $\Omega_\text{tor}$ plays the role of the lapse function in GR. The
gravitational potential $\phi_g$ enters via:
\[
  \Omega_\text{tor}(\vec{x}) = \omega_c\left(1 - \frac{\phi_g(\vec{x})}{c^2}\right),
\]
recovering gravitational time dilation to leading order. A derivation of the full
Poisson equation from the incompressible fluid Euler equations is given in
SST-Canon \cite{Iskandarani2025Canon}.

% ============================================================
\section{Falsifiable Predictions}\label{sec:falsifiers}
% ============================================================

SST makes the following falsifiable predictions beyond mass values:

\begin{enumerate}
  \item \textbf{Knot-mass correspondence.} The masses of the $5_2$ and $6_1$ knots
        fix the quark mass ratio. SST predicts
        $m_d / m_u = V_{6_1}/V_{5_2} = 3.16396/2.82812 \approx 1.119$.
        Lattice QCD gives $m_d/m_u \approx 1.9$--$2.2$; the current SST value is a
        factor $\sim 2$ low. This discrepancy is a falsifier if it cannot be resolved
        by the inter-knot coupling corrections.

  \item \textbf{Swirl speed as universal constant.} The prediction
        $v_{\circlearrowleft} = c\alpha/2$ should appear as a limiting surface
        acoustic velocity in any solid-state vortex system. Independent experimental
        confirmation from SAW/MEMS devices \cite{Iskandarani2026AIP} is consistent
        at $< 0.1\%$.

  \item \textbf{Achiral knots decouple.} Knots with zero amphichirality
        (e.g.\ figure-eight $4_1$) should carry no electromagnetic coupling.
        This predicts a gravitationally interacting but electromagnetically inert
        dark sector. Its density is constrained by cosmological observations.

  \item \textbf{NLS discreteness in BEC vortices.} The delay-induced mass ladder
        should appear as anomalous discrete frequency shifts in Bose-Einstein
        condensate vortex rings under parametric driving at $f_c$.

  \item \textbf{Neutron mass correction.} Once the inter-knot gear coupling
        coefficient is computed analytically from knot linking numbers, the
        neutron mass prediction should improve to $< 1\%$. Failure to do so
        would challenge the knot-assignment for the down quark.
\end{enumerate}

% ============================================================
\section{Discussion}
% ============================================================

The central result of this paper is the identification:
\[
  \tau_\text{vortex} = \frac{1}{f_c},
\]
which bridges classical fluid delay dynamics and quantum Compton-scale physics
through the geometry of the $R = 2r$ torus. Combined with the cavitation-limit
anchoring of $F_{\circlearrowleft}^{\max}$ to the gravitational maximum-force
invariant, this eliminates the circular dependence on electron mass that plagued
earlier VÆM formulations.

The framework unifies three aspects of physics that are usually treated separately:
\begin{itemize}
  \item \textbf{Mass} as topologically trapped rotational energy
  \item \textbf{Electromagnetism} as poloidal vortex coupling
  \item \textbf{Gravity} as toroidal vortex clock rate, i.e.\ the local rate of
        the delay feedback loop
\end{itemize}
All three emerge from a single incompressible fluid with two scale parameters:
$r_c$ (geometric) and $v_{\circlearrowleft}$ (kinematic), both derivable from $\alpha$,
$\hbar$, $c$ without mass input.

The NLS extension of the delay paper \cite{Iskandarani2026AIP} provides the field
equation that governs vortex mode selection. The quasi-geostrophic stream-function
equation identified in 2014 \cite{Iskandarani2014VAM} is recovered as the
irrotational limit of this NLS in a rotating frame --- confirming that the 2014
intuition was correct in its essential structure.

% ============================================================
\section{Conclusion}
% ============================================================

We have presented SST-70, extending the delay-induced mode-selection mechanism of
Iskandarani (2026) to the full nonlinear Schr\"{o}dinger framework of vortex swirl
clocks. The key results are:

\begin{enumerate}
  \item The Compton frequency is the vortex self-resonance frequency:
        $f_c = v_{\circlearrowleft}/(2\pi r_c)$.
  \item All SST parameters are anchored to $(c, \hbar, \alpha, G)$ without
        circularity.
  \item Knot topology provides the mode selector; discrete stable masses emerge
        from stability constraints on the delay differential equation.
  \item Mass predictions agree with CODATA to within $3\%$ for baryons and
        $< 10^{-9}\%$ for the electron, using no particle mass as input.
\end{enumerate}

The neutron discrepancy and the quark mass ratio remain open problems, providing
clear targets for the next iteration of SST.

% ============================================================
\appendix
\section{SST Canonical Constants}
% ============================================================

\begin{table}[H]
\centering
\begin{tabular}{llll}
\toprule
\textbf{Symbol} & \textbf{Value} & \textbf{Unit} & \textbf{Origin} \\
\midrule
$v_{\circlearrowleft}$ & $1.09384563\times10^6$ & m s$^{-1}$ & $c\alpha/2$ \\
$r_c$           & $1.40897017\times10^{-15}$ & m & $\lambda_c\alpha/4\pi$ \\
$\rho_c$        & $3.89344\times10^{18}$ & kg m$^{-3}$ & $4F_{\circlearrowleft}^{\max}/(\pi\alpha^2 c^2 r_c^2)$ \\
$F_{\circlearrowleft}^{\max}$ & $29.0535$ & N & $e^2/(16\pi\varepsilon_0 r_c^2)$ \\
$\varphi$       & $1.61803399$ & --- & $(1+\sqrt{5})/2$ \\
$\alpha$        & $7.2973526\times10^{-3}$ & --- & CODATA \\
\bottomrule
\end{tabular}
\caption{SST primary constants. All derivable from $(c,\hbar,e,G)$.}
\end{table}

% ============================================================
\appendix
\section{Knot Volume Table}
% ============================================================

\begin{table}[H]
\centering
\begin{tabular}{llll}
\toprule
\textbf{Knot} & \textbf{Hyp.\ Volume} & \textbf{Chirality} & \textbf{SST assignment} \\
\midrule
$3_1$ trefoil      & torus knot  & Yes & Electron \\
$4_1$ figure-eight & 2.0298      & No  & Dark sector \\
$5_1$ cinquefoil   & torus knot  & Yes & Higher lepton \\
$5_2$ twist        & 2.82812     & Yes & Up quark \\
$6_1$ Stevedore    & 3.16396     & Yes & Down quark \\
\bottomrule
\end{tabular}
\caption{Knot assignments in SST. Achiral knots decouple from the mass mechanism.}
\end{table}

% ============================================================
\begin{thebibliography}{99}

\bibitem{Iskandarani2013}
O.\ Iskandarani,
\textit{Een Hypothese over de Constantes in de Kwantum Mechanica,
Verklaard met de Klassieke Logica}
[A Hypothesis on the Constants of Quantum Mechanics, Explained with
Classical Logic],
Public Google Document, Groningen, July 2013.\\
\url{https://docs.google.com/document/d/1lEji7B5K_CmBNnLKiJ-ngB7Q_1Gi9_qDPKKPewGCzyU/edit}\\
\textit{(Establishes $C_e$, $F_{\max}$, $h = 4\pi F_{\max}r_c^2/C_e$,
$\alpha = 2C_e/c$; priority document.)}

\bibitem{Iskandarani2014a}
O.\ Iskandarani,
\textit{Korte Fundamentele Samenvatting van de Eigenschappen van de \AE ther}
[Short Fundamental Summary of the Properties of the \AE ther],
Public Google Document, Groningen (2014).\\
\url{https://docs.google.com/document/d/15J9kZOfHtrz_9q2-Q8jThsL3oJfuhTHg9qk_6P41MNc/edit}\\
\textit{(Contains: quark formation from vortex fusion; three-phase aether model;
the statement ``$c$ is not the speed of light but the ratio of electrostatic
to electromagnetic charge units'' --- anticipating $v_\circlearrowleft = c\alpha/2$.)}

\bibitem{Iskandarani2014VAM}
O.\ Iskandarani,
\textit{Samenvatting van de Onsamendrukbare Niet-Viskeuze Vloeistof \AE ther}
[Summary of the Incompressible Inviscid Fluid Æther],
Public Google Document, Groningen, 2014.\\
\url{https://docs.google.com/document/d/1dDfA-82VVl-S6ZuFCwhlze1tlowB9FZnt2-VPTWFqPM/edit}\\
\textit{(Establishes: Maxwell field lines = Helmholtz vortex lines = Einstein
geodesics; time as kernel vortex circulation $d\tau = d\Gamma_\text{kernel}$.)}

\bibitem{Iskandarani2015Notes}
O.\ Iskandarani,
\textit{Rotating Frame Fluid Mechanics and the Structure of Quantum Mechanics},
Unpublished working notes, Groningen (2015).\\
\textit{(Identifies Coriolis term $2\vec{\omega}\times\vec{U}$ with the
magnetic force; notes that QM implicitly combines gravitational and centrifugal
potentials.)}

\bibitem{Iskandarani2025Canon}
O.\ Iskandarani,
\textit{Swirl-String Theory: Canon v0.7.8},
Zenodo preprint (2025).
\url{https://zenodo.org/records/18913147}

\bibitem{Iskandarani2026AIP}
O.\ Iskandarani,
``Delay-Induced Mode Selection in Circulating Feedback Systems,''
\textit{AIP Advances} (2026, in review, MS\#2064597).\\
\textit{(Proves delay alone is sufficient for spectral discreteness in
closed circulation loops; foundational for the present work.)}

\bibitem{KnotInfo}
C.\ Livingston and A.\ H.\ Moore,
\textit{KnotInfo: Table of Knot Invariants},
\url{https://knotinfo.math.indiana.edu}, (2023).

\bibitem{BurdeZieschang}
G.\ Burde and H.\ Zieschang,
\textit{Knots}, 2nd ed., de Gruyter, Berlin (2003).

\bibitem{Willis2012}
J.\ R.\ Willis,
``The construction of effective relations for waves in a composite,''
\textit{Comptes Rendus M\'ecanique} \textbf{340}, 181--192 (2012).

\bibitem{Kleckner2013}
D.\ Kleckner and W.\ T.\ M.\ Irvine,
``Creation and dynamics of knotted vortices,''
\textit{Nature Physics} \textbf{9}, 253--258 (2013).

\bibitem{Helmholtz1858}
H.\ von Helmholtz,
``{\"U}ber Integrale der hydrodynamischen Gleichungen, welche den
Wirbelbewegungen entsprechen,''
\textit{Journal f{\"u}r die reine und angewandte Mathematik} \textbf{55},
25--55 (1858).

\bibitem{Kelvin1867}
W.\ Thomson (Lord Kelvin),
``On Vortex Atoms,''
\textit{Philosophical Magazine} \textbf{34}, 15--24 (1867).

\bibitem{Maxwell1861}
J.\ C.\ Maxwell,
``On Physical Lines of Force,''
\textit{Philosophical Magazine} \textbf{21} (1861).

\end{thebibliography}

\end{document}
